% page 140 du fac-simile (image p-139 du scan)
% bloc 152.4 x 233.3 mm ; marge gauche 8.0 mm
\pgc{-12.700mm}{0.000mm}{
\l{\cel{3}130}
\l{}
\l{\sou{elektroskopo}. Aparato per qua onu povas konstatar la elek-}
\l{\cel{3}trozeso. - DEFIRS.}
\l{}
\l{\sou{elektroskopio}. Studiado di la elektroskopi e di la apliko-}
\l{\cel{3}kazi elektroskopala. - DEFIRS.}
\l{}
\l{\sou{elektuario}. (\sou{Medic.}) Medikamento koheranta mola, ek subs-}
\l{\cel{3}tanci diversa (pulvero, pulpi) interligita per siropo}
\l{\cel{3}o vino. - DEFIRS.}
\l{}
\l{\sou{elemento}. I. (\sou{che la antiqui}). Parto konstitucanta di kozo.}
\l{\cel{3}- II. (\sou{che la moderni}). Korpo quan onu kredas simpla,}
\l{\cel{3}kemiale. - III. En la linguo ordinara : singla ek la}
\l{\cel{3}kozi di qui la asemblo e la kombino formacas kozo altra.}
\l{\cel{3}- IV. (\sou{pri cienco}). La principi prima sur qui lu fond-}
\l{\cel{3}esas, e per qui onu komencas lernar lu. - DEFIRS.}
\l{}
\l{\sou{elevar}. -trans.) Transportar a nivelo supera. (\sou{anke meta}-}
\l{\cel{3}\sou{fore}).\cel{2}EFIRS.}
\l{}
\l{\sou{elevaciono}. Desegnuro di la formi extera di ula korpo vidata}
\l{\cel{3}en direciono horizontala. - EFS.}
\l{}
\l{\sou{elevatoro}. I. (\sou{anat.}) Muskulo qua levas ulo. - II. (\sou{mek}a-}
\l{\cel{3}\sou{niko}). Aparato uzata por levar la navi en la reparo-doki.}
\l{\cel{3}- DEFIRS.}
\l{}
\l{\sou{elfo}. (\sou{mitol. Skandinava}). Koboldo di atmosfero.- DEFIRS.}
\l{}
\l{\sou{eliminar}.\cel{2}(\sou{trans.}) (\sou{tekniko, matem.}) Ekirigar, eskartar.}
\l{\cel{3}DEFIS.}
\l{}
\l{\sou{elipso}. I. (\sou{gramatiko}). Figuro di retoriko qua konsistas ye}
\l{\cel{3}supresar un o plura vorti quin la lektero od askoltero}
\l{\cel{3}devos suplear.- II. (\sou{geom}.) Kurvo qua rezultas de la seko}
\l{\cel{3}di kono rekta per plano obliqua ye la axo e ye omna}
\l{\cel{3}aristi. - DEFIRS.}
\l{}
\l{elipsoida. Elipso-forma. - DEF.}
\l{}
\l{\sou{elito}. To quo selektesas ek korpo kom lo maxim bona.-DEF.}
\l{}
\l{elitro.\cel{2}(\sou{zool.}) Alo rezistiva qua kovras, che ula insekti,}
\l{\cel{3}ali plu dina e plu tenua. - DEF.}
\l{}
\l{\sou{elixiro}. Preparajo farmaciala, qua rezultas de la mixo de}
\l{\cel{3}ula siropi kun alkoholaji. - DEFIRS.}
\l{}
\l{\sou{elizeo}. I. (\sou{che la antiqui}) Regioni di inferno ube, segun}
\l{\cel{3}la antiqui, la homi vertuoza, la heroi, sejornis pos}
\l{\cel{3}lia\cel{2}orto ed ube printempo esis eterna. - II. (\sou{metaf}.)}
\l{\cel{3}Loko delicoza. - DEFIRS.}
}
